\documentclass[12pt]{article}
\usepackage[utf8]{inputenc}
\usepackage[T1]{fontenc}
\usepackage{hyperref}
\title{Urali Language Resources for K-12 Curriculum Development}
\author{Compiled by Claude Sonnet 5 (claude-sonnet-5) for Andr\'as Kornai}
\date{}
\begin{document}
\maketitle
\section{Introduction}
Urali (ISO 639-3 url, Glottocode ural1274) is a Dravidian language of southern India, described as dormant. Sources disagree on its classification. Wikipedia treats it as a dialect of Irula, while other sources equate ``Urali Kurumba'' with Betta Kurumba. The 2011 census figure found is 12,986, and the speakers' ethnic identity should not be assumed to match the language. No Urali school materials turned up in search.
\section{Primary Educational Resources}
None found. No primers, textbooks or school programmes for Urali appeared in search results.
\section{Linguistic and Reference Materials}
\begin{itemize}
\item Urali language (Wikipedia): \url{https://en.wikipedia.org/wiki/Urali_language}
\item Kurumba languages (Wikipedia), which lists Urali Kurumba as an alias of Betta Kurumba: \url{https://en.wikipedia.org/wiki/Kurumba_languages}
\item Glottolog entry for Betta Kurumba (Urali Kurumba): \url{https://glottolog.org/resource/languoid/iso/xub}
\item Betta Kurumba bibliography, Endangered Languages Project. It lists Gail Coelho's 2003 grammar thesis (University of Texas at Austin), her 2018 \emph{Annotated Texts in Betta Kurumba} with a grammatical sketch and glossary, and a 2005 HRELP documentation project: \url{https://www.endangeredlanguages.com/lang/1453/bibliography}
\item De Gruyter record for the 2018 Betta Kurumba texts volume: \url{https://www.degruyterbrill.com/document/isbn/9789004378247/html}
\end{itemize}
These concern Betta Kurumba, which may or may not be the same variety as Urali. No dictionary was found.
\section{Religious and Community Resources}
None found.
\section{Digital and Audio Resources}
None found for Urali specifically.
\section{Curriculum Development Context}
Little exists for Urali. Any curriculum work would first need to settle which variety the term refers to (Irula dialect or Betta Kurumba), then draw on the Betta Kurumba documentation and on community consultation. The language is not known to be taught in schools.
\end{document}